\subsection*{Quantitative azimuthal moment restoration at (A.11)}
\label{subsec:exact-A11}
The source's angular correction must be evaluated with the prescribed
translated copies of $\sigma'$, rather than arbitrary moment bumps.
We give an explicit inverse and a quantitative nonlinear correction
on the post-pulse constant-slope interval. This analysis concerns
only the two moment equalities in (A.11); it does not freeze the
preceding stress-cone parameters or the later exterior matching time.

Put $y=\log(X/X_1)$, where $X_1$ is the center of the first
correction bump. The second center is $y=2$, and the unedited
profile has $\ell=-\lambda$ between them. Let
\[
 \beta(t)=\frac{10}{3}\sigma'\!\left(\frac{10t}{3}+\frac12\right),
 \quad \beta_1(y)=\beta(y),\quad\beta_2(y)=\beta(y-2),
 \quad \operatorname{supp}\beta\subset[-3/20,3/20].
\]
Here $\int\beta=1$ and, by the corrected source-bump lemma,
$\int\beta^2\le640/3$. Write
$E_c=E_{\rm base}(1+c_1\beta_1+c_2\beta_2)$.
The source's angular-moment equation and its pressure-neutrality
equation are precisely
\begin{equation}\label{eq:A11-two-exact}
 \begin{split}
 u(c_1+xc_2)&=\Delta_I,\\
 2v(c_1+zc_2)+w(c_1^2+zc_2^2)&=0,
 \end{split}
\end{equation}
where the positive source-bump integrals are
\begin{equation}\label{eq:A11-weights}
\begin{gathered}
 u=\int e^{(1-\lambda)t}\beta(t)\,dt,\qquad
 v=\int e^{-(1+2\lambda)t}\beta(t)\,dt,\qquad
 w=\int e^{-(1+2\lambda)t}\beta(t)^2\,dt,\\
 x=e^{2(1-\lambda)},\qquad z=e^{-2(1+2\lambda)},\qquad
 \Delta_I=\frac1{1-\lambda}-\frac{I(X_1)}{X_1H(X_1)}.
\end{gathered}
\end{equation}
Equivalently, put $\alpha=\Delta_I/u$ and eliminate
$c_1=\alpha-xc_2$. Then the \emph{actual} two-equation
correction is reduced to the scalar quadratic
\[
 \mathsf A c_2^2+\mathsf Bc_2+\mathsf C=0,\quad
 \mathsf A=w(x^2+z),\quad
 \mathsf B=2v(z-x)-2w\alpha x,\quad
 \mathsf C=2v\alpha+w\alpha^2.
\]
On the small branch supplied below $\mathsf B<0$ and the
discriminant is positive. A cancellation-safe evaluation is
$c_2=2\mathsf C/[\sqrt{\mathsf B^2-4\mathsf A\mathsf C}-\mathsf B]$,
followed by $c_1=\alpha-xc_2$. This is an exact specification
of the two source correction coefficients once $\Delta_I$ is known.

The first identity follows by evaluating the change in
$I=\int H\,dX$ and normalizing by $X_1H(X_1)$, including the
fact that $I/(XH)$ approaches $1/(1-\lambda)$ under the
unedited constant-slope flow. The second is exactly
$\int(E_c^2-E_{\rm base}^2)\,dy=0$: disjoint supports
eliminate its mixed quadratic term.

\begin{proposition}[Uniform two-bump inverse and quadratic bound]
\label{prop:A11-matrix-bound}
For $0<\lambda\le1/20$, the Jacobian and quadratic part of
\eqref{eq:A11-two-exact} satisfy
\[
 B_{11}=\begin{pmatrix}u&ux\\2v&2vz\end{pmatrix},
 \qquad Q(c,c)=\binom{0}{w(c_1^2+zc_2^2)},
 \qquad\|B_{11}^{-1}\|_\infty<\frac32,\quad\|Q\|<320.
\]
In particular, its determinant is bounded away from zero
uniformly as $\lambda\downarrow0$.
\end{proposition}
\begin{proof}
Since $|t|\le3/20$ on the support, both $u$ and $v$ belong
to $[4/5,5/4]$ by $e^{-1/5}\ge4/5$ and
$e^{1/5}\le5/4$. Also $6<x<8$ and $0<z<1/7$.
The determinant is $2uv(z-x)$, with magnitude exceeding
$2(4/5)^2(6-1/7)=1312/175$. Row sums of its inverse
are bounded by
\[
 \frac{2(5/4)(1/7)+(5/4)8}{1312/175}
 =\frac{3625}{2624}<\frac32,
 \qquad
 \frac{2(5/4)+5/4}{1312/175}
 =\frac{2625}{5248}<\frac32.
\]
The bound on the prescribed bump implies
$w\le(5/4)(640/3)$, and hence
$\|Q\|\le w(1+z)<6400/21<320$.
\end{proof}

\begin{proposition}[Source moment discrepancy at the first bump]
\label{prop:A11-discrepancy}
Suppose the prescribed pre-pulse and interpolation stages
satisfy their exact source bounds, namely
$\ell\ge-\lambda-1/10$ and
$\ell=-\lambda$ throughout the following
$30\log(1/\lambda)$ interval. Then
\begin{equation}\label{eq:A11-delta-small}
 |\Delta_I|<64\lambda^{28}.
\end{equation}
\end{proposition}
\begin{proof}
For $r=I/(XH)$, the exact source transport law is
$r'+(1+\ell)r=1$. At the reference radius,
$r=5/8$, and during all pre-bump stages
$1+\ell\ge17/20$, so $0\le r\le20/17$.
Under $\ell=-\lambda$, the difference from the fixed
point $1/(1-\lambda)$ decreases at rate $1-\lambda$.
The first bump lies three units before the end of a plateau
of length $30\log(1/\lambda)$. Consequently
\[
 |\Delta_I|<3e^{3(1-\lambda)}\lambda^{30(1-\lambda)}
 \le3e^3\lambda^{57/2}<64\lambda^{28},
\]
since $e^3<21$ and $0<\lambda\le1/20$.
\end{proof}

\begin{theorem}[The two actual (A.11) corrections preserve negative slope]
\label{thm:A11-true-corrections}
For $0<\lambda\le1/20$ the source's two equations
\eqref{eq:A11-two-exact} admit a unique solution in the
small-correction ball with
\begin{equation}\label{eq:A11-c-coeff}
 \|c\|_\infty\le192\lambda^{28}.
\end{equation}
The edited azimuthal field is positive, and on both bump
supports the actual logarithmic slope satisfies
\begin{equation}\label{eq:A11-slope}
 -\frac{3\lambda}{2}\le\ell_c\le-\frac\lambda2.
\end{equation}
Off the supports, $\ell_c=-\lambda$.
\end{theorem}
\begin{proof}
Propositions~\ref{prop:A11-matrix-bound} and
\ref{prop:A11-discrepancy} give
$8\|B_{11}^{-1}\|^2\|Q\||\Delta_I|
 <368640\lambda^{28}<1$. The contraction lemma A.2
of the source therefore supplies a unique small solution
with $\|c\|_\infty\le2(3/2)|\Delta_I|
\le192\lambda^{28}$.
For completeness, the actual step $\sigma$ has derivative
bounds $\|\sigma'\|_\infty\le64$ and
$\|\sigma''\|_\infty\le17152$.
On $[1/4,3/4]$, write
$g=-t^{-2}+(1-t)^{-2}$ and
$\sigma=(1+e^{-g})^{-1}$; then
$|g'|\le256$, $|g''|\le3072$, giving the asserted
second-derivative bound. For $0<t\le1/4$, put $s=1/t\ge4$.
The tail derivatives satisfy
$|\sigma'|\le3s^3e^{-s^2+16/9}<1$ and
$|\sigma''|\le7s^6e^{-s^2+16/9}<1$;
symmetry gives the remaining quarter.
Thus $\|\beta\|_\infty\le640/3$ and
$\|\beta'\|_\infty<200000$.
The two supports are disjoint, so
$\|\sum_i c_i\beta_i\|_\infty
\le(640/3)192\lambda^{28}<1/2$.
Moreover
\[
 |\ell_c+\lambda|
 =\left|\frac{\sum_i c_i\beta_i'}{1+\sum_i c_i\beta_i}\right|
 <76800000\lambda^{28}<\frac\lambda2.
\]
This proves positivity and the slope interval.
\end{proof}

\paragraph{Explicit interpolation and terminal slope parameters.}
One may fix $T_f=512$ in (A.10): since
$0\le J_0(\eta)=\log(1+\eta^2)\le\log2<7/10$
and $\|\sigma'\|_\infty\le64$, its slope satisfies
\begin{equation}\label{eq:A11-Tf-fixed}
 -\lambda-\frac7{80}\le\ell_{\rm interp}\le-\lambda.
\end{equation}
Likewise the terminal cutoff coefficient in (A.12) may be
fixed as $c_o=1/256$, which yields
$0\le f_o'/f_o<h/4$ for $0<h\le1/100$.
For every $T_d>10$, the explicit parameter expressions
\begin{equation}\label{eq:A11-partial-tuple}
 T_f=512,\quad c_o=1/256,\quad
 \lambda=\frac{10^{-6}}{T_d+513},\quad
 h=\frac{\lambda e^{-T_d}}{10^6}
\end{equation}
satisfy all smallness assumptions used in this section and
in the preceding forward/backward pulse-energy inequalities.
This is a \emph{partial} quantitative parameter selection:
it does not certify an admissible choice of $M_d,P_*$,
all later stress-cone constants, or the unique stopping length
of the terminal $Q_s$ matching interval. No full profile,
Appendix-B contraction, unforced similarity corrections or
physical Navier--Stokes solution is thereby established.
